\section{Special cover strata and the old axis summand}

\subsection{C44-C01: the old variable classes do not disappear}
\begin{proposition}[C44-C01]\label{prop:old-axis-cover}
Use JAC1964 H9-T18. Write $A$ for the abelian fourfold, $B$ for its
surface axis, $X_0=\operatorname{Bl}_B A$, $M_0=A\times\mathbf P^1$,
and $k_0:X_0\hookrightarrow M_0$ for the incidence embedding.
For a finite separable cover $\pi:C\to\mathbf P^1$ of degree $b$,
unramified over the critical values, put
$X_C=X_0\times_{\mathbf P^1}C$, $M_C=A\times C$, and let $g:X_C\to X_0$
and $k_C:X_C\hookrightarrow M_C$ be the maps. On algebraic middle classes,
\[
 \Alg^2(X_C)\cap\ker k_{C*}
 =g^*(\Alg^2(X_0)\cap\ker k_{0*})\ \mathbin\perp\
   (\Alg^2(X_C)\cap\ker k_{C*}\cap\ker g_*).
 \tag{C44.C01.1}
\]
The first summand is canonically the old axis space
$A_B=\ker(\kappa_*:\Alg^1(B)\to H^6(A)(1))$ and has form
\[
          (w,w')\longmapsto-b\int_Bww'.                 \tag{C44.C01.2}
\]
It is positive definite, by the surface Hodge index theorem, without 032.
Consequently the remaining H9-T18 sign is exactly the sign on the second
summand in (C44.C01.1). Statements that generic covers add no algebraic
middle classes refer to this new part (or the full trace kernel), and do
not say that the entire H9-T18 variable space vanishes.
\end{proposition}
\begin{proof}
The finite flat map $g$ has $g_*g^*=b$. Projection formula makes
$P=b^{-1}g^*g_*$ a self-adjoint algebraic idempotent. The Cartesian
incidence square gives
$k_{C*}g^*=(1_A\times\pi)^*k_{0*}$, and proper functoriality gives
$k_{0*}g_*=(1_A\times\pi)_*k_{C*}$. Thus $P$ preserves
$\ker k_{C*}$, and its image there is exactly the first summand of
(C44.C01.1); its kernel is the displayed second summand. Both are
algebraic subspaces, since $P$ and $1-P$ are algebraic. Self-adjointness
proves orthogonality.

For the old space use the codimension-two blowup decomposition from
H9-T7. A middle class has unique form $q^*z+e_*p^*w$ with
$z\in H^4(A)$ and $w\in H^2(B)$. With
$t=[\mathrm{point}]\in H^2(\mathbf P^1)$, incidence pushforward is
\[
 k_{0*}(q^*z+e_*p^*w)=(hz+\kappa_*w)+zt.
\]
The exceptional divisor is $B\times\mathbf P^1$ as a variety over
$A\times\mathbf P^1$; this proves the second term's formula, while
the first follows from $[X_0]=h+t$. Kunneth independence forces $z=0$
and $\kappa_*w=0$ in the kernel. This is the asserted axis identification.
H9-T7 equation (H9.24) gives its old form $-\int_Bww'$. Pullback under
$g$ multiplies intersections by $b$, proving (C44.C01.2). Its strict
positivity is precisely the surface input used in H9-T8. No sign input
on $X_C$ has been used.
\end{proof}

\paragraph{Archive and predecessor implications.}
This sharpens the wording of the cover gate for cote 44, Th.~5,
scan 004 and JAC1964 QH9-cover-primitivep. It uses H9-T7, H9-T8,
H9-T18 and finite-flat push-pull; it does not re-prove their height
identification. Cote 16, scan 050, language is algebraic Gysin and
restriction with their correct Lefschetz weights. In Chow motives,
$b^{-1}g^*g_*$ is the old-part projector; the unresolved piece is its
complement inside the primitive middle projector.

\subsection{C44-C02: every nonextremal fixed splitting stratum of triple covers}
\begin{theorem}[C44-C02]\label{thm:maroni}
Fix the abelian pencil of H9-T23 and integers
\[
                  1\le e\le f<2e,
                  \qquad g=e+f-2\ge2.
\]
In every characteristic, let $U_{e,f}$ be the parameter space of smooth
connected triple covers of the fixed target $\mathbf P^1$, unramified
above the pencil's critical set, with framed Tschirnhausen bundle
$\mathcal O(e)\oplus\mathcal O(f)$. It is a nonempty irreducible open
subset of the space of twisted binary cubics. At its geometric generic
point, every algebraic class in the new middle trace kernel of the
pullback pencil is zero. The exact norm-height formula and radical are
\[
 h_L(z,w)=\tfrac13h_K(\operatorname{cor}z,\operatorname{cor}w),
 \qquad \operatorname{rad}h_L=
       \operatorname{cor}^{-1}(\operatorname{rad}h_K).       \tag{C44.C02.1}
\]
For abelian-fourfold pencils the full homological height is therefore
semipositive, conditionally on 032/Milne as in H9-T8. The vanishing
assertion and norm identity themselves do not use 032.

For $d=f-e>0$ the unframed stratum has codimension $d-1$ in the full
fixed-target smooth triple-cover stack. Thus this treats in particular
the first unbalanced stratum for even $g\ge6$ (a divisor), and for odd
$g\ge9$ (codimension two). It does not assert anything for every closed
point of these strata or for the extremal splitting $f=2e$.
\end{theorem}
\begin{proof}
For $E=\mathcal O(e)\oplus\mathcal O(f)$, Wood's cubic-algebra
equivalence, in the exact form already used in H9-T24 equations
(H9.136)--(H9.139), gives the affine coefficient space
\[
 P_E=H^0\bigl(\mathcal O(2e-f)\oplus\mathcal O(e)
                \oplus\mathcal O(f)\oplus\mathcal O(2f-e)\bigr).
\]
All degrees are nonnegative, so $\dim P_E=2g+8$. The smooth,
geometrically connected, and specified unramified conditions define an
open $U_E$. Each point has the required quotient bundle by the cubic
equivalence. It remains to show nonemptiness and the monodromy property.

Set
\[
             a=2e-f-1\ge0,\qquad c=f-e\ge0.
\]
Then $2a+3c=g$, $a+c+1=e$, and $a+2c+1=f$. Apply H9-T23 with
these particular parameters (its proof allows arbitrary nonnegative
$a,c$, not only the balanced choices used later in H9-T24).
In characteristic different from three, the Kummer normalization has
the relative eigensheaves
$\mathcal O\oplus\mathcal L^{-1}\oplus\mathcal L^{-2}(D_2)$,
where $\deg\mathcal L=e$ and $\deg D_2=a+1$. Its quotient dual is
$\mathcal O(e)\oplus\mathcal O(2e-a-1)=E$.
All exponent-two points and the line bundle are fixed over the moving
branch parameter, so this is an isomorphism of relative bundles.
In characteristic three, there are $a+1$ simple poles and $c$ double
poles. The exact normalized basis calculation (H9.144) gives quotient
dual degrees $a+c+1=e$ and $a+2c+1=f$, again with constant pole
divisors. After choosing a framing, in each case this constructs an
actual morphism from the H9-T23 parameter curve into $U_E$.

The universal pullback pencil is smooth proper, and
$1-\frac13Q^*Q_*$ defines a lisse middle-cohomology direct summand,
exactly as in H9-T24 equations (H9.146)--(H9.147). Its restriction to
the constructed curve has no nonzero finite-orbit vectors by H9-T23's
middle-convolution or Fourier calculation. A finite-orbit vector for
the fundamental group of $U_E$ would remain such after restriction to
this curve. Hence none exists on $U_E$. Every algebraic class at the
geometric generic point is defined over a finite extension of the
parameter function field and consequently has a finite monodromy
orbit. Its new projection must vanish. The actual-cycle push-pull and
admissibility argument in H9-T24 then gives (C44.C02.1) and the radical.
This uses precisely the old result, with the newly specified parameter
stratum, rather than a generic conclusion specialized to a closed point.

For codimension, $\dim\operatorname{Aut}(E)=d+3$ if $d>0$ and is
$4$ if $d=0$: compute $H^0(\End E)$ from its four line summands.
Thus the $d>0$ stratum has dimension $2g+5-d$, while the balanced
stratum has dimension $2g+4$. Smooth separable covers have finite
automorphism groups, so the quotient-stack dimension calculation is
valid also in characteristics two and three. The difference is $d-1$.
The first nonbalanced value is $d=2$ for even $g$ and $d=3$ for odd
$g$; $f<2e$ gives exactly the genus bounds stated above.
\end{proof}

\paragraph{Sources and scope.}
Wood, \emph{Parametrizing quartic algebras over an arbitrary base},
Theorem 2.1 and equation (1), printed pp.~1073--1074; JAC1964 H9-T23,
(H9.127)--(H9.132), and H9-T24, (H9.136)--(H9.147). These are the
precise predecessor inputs. The new conclusion concerns geometric
generic points of special Maroni strata, over algebraic closures of
their parameter fields. It is not a proof for all special covers over
$\overline{\F}_p$. Archive target: cote 44, scan 004, Th.~5 and its
fibration margin; predecessor row QH9-cover-primitivep is narrowed in
this specified stratum scope, not closed.
